\documentclass[11pt]{article}
\usepackage[T1]{fontenc}
\usepackage[utf8]{inputenc}
\usepackage{lmodern,microtype}
\usepackage[a4paper,margin=27mm]{geometry}
\usepackage{amsmath,amssymb,amsthm,mathtools}
\usepackage{booktabs,array,longtable,xcolor,listings}
\usepackage{fancyhdr}
\usepackage[unicode,hidelinks]{hyperref}
\hypersetup{pdftitle={Period Obstructions for Weighted Binary Expansions},pdfauthor={Keio},pdfsubject={Denominator valuations, periodic expansions, and the Borwein--Loring periodicity question}}
\newtheorem{theorem}{Theorem}[section]
\newtheorem{lemma}[theorem]{Lemma}
\newtheorem{corollary}[theorem]{Corollary}
\newtheorem{proposition}[theorem]{Proposition}
\theoremstyle{definition}
\newtheorem{definition}[theorem]{Definition}
\newtheorem{example}[theorem]{Example}
\theoremstyle{remark}
\newtheorem{remark}[theorem]{Remark}
\DeclareMathOperator{\ord}{ord}
\setlength{\emergencystretch}{2em}
\setlength{\headheight}{14pt}
\pagestyle{fancy}
\fancyhf{}
\fancyhead[L]{\small Period obstructions for weighted binary expansions}
\fancyhead[R]{\small 2 November 2026}
\fancyfoot[C]{\thepage}
\lstset{language=Python,basicstyle=\ttfamily\scriptsize,
  keywordstyle=\color{blue!45!black},commentstyle=\color{black!55},
  stringstyle=\color{black!75},showstringspaces=false,
  columns=fullflexible,keepspaces=true,breaklines=true,
  frame=single,rulecolor=\color{black!20},framesep=5pt,
  xleftmargin=6pt,xrightmargin=6pt,aboveskip=10pt,belowskip=10pt}
\title{\textbf{Period Obstructions for\\Weighted Binary Expansions}\\[5pt]
\large Denominator valuations and a 1990 periodicity question}
\author{Keio}
\date{2 November 2026}
\begin{document}
\maketitle
\begin{abstract}
For binary digits $\varepsilon_n$, consider the weighted value
$X=\sum_{n\ge1}n\varepsilon_n/2^n$. We prove that an eventually periodic
representation with a nonconstant tail forces an odd prime square in the
reduced denominator of $X$. Consequently, no rational with a nontrivial
squarefree odd denominator part has any eventually periodic representation;
in particular this holds for $1/3$. This rules out the universal periodicity
possibility raised by Borwein and Loring in 1990. More precisely, an exact
valuation identity and a finite closure under prime divisors of
multiplicative orders give an explicit integer $D_Q$, depending only on
the odd denominator $Q$ of $X$, such that every ordinary periodic-tail
denominator divides $D_Q$ and every minimal period divides
$\ord_{D_Q}(2)$. We give full proofs, a general-base extension, boundary
examples, and an exact reproducible verifier. This does not resolve
Erd\H{o}s Problem~\#261; historical priority of the additional results is
not certified by the bounded literature search reported here.
\end{abstract}

\section{The historical question and the auxiliary problem}
Borwein and Loring call an expression
\begin{equation}\label{eq:weighted-def}
 x=\sum_{n\ge1}\frac{n\varepsilon_n}{2^n},\qquad
 \varepsilon_n\in\{0,1\},
\end{equation}
a $*$-binary representation. On printed page~391 of~\cite{BL1990}, they
leave open the possibility that every rational in the range admits a
periodic representation and single out $1/3$. Their discussion explicitly
allows a period to start after a finite initial segment. The question
therefore concerns eventual periodicity, which is the convention used
throughout this manuscript.

The quantifier is essential: ruling out one algorithm's periodic output
would not rule out other representations. Corollary~\ref{cor:squarefree-periodicity}
rules out \emph{every} eventually periodic representation of $1/3$, with
an arbitrarily long preperiod allowed. Thus it gives a negative answer
to that particular universal possibility without changing its meaning.
It is an unnumbered question in the cited paper, rather than an
established named conjecture or the whole of Erd\H{o}s~\#261.

Our auxiliary problem is to find a finite, explicit bound for the
ordinary binary tail denominators that can occur when the \emph{weighted}
value has a prescribed odd denominator $Q$. A periodic ordinary binary value has denominator
dividing $2^r-1$, but differentiating the corresponding generating
function introduces a second power of this factor. The key question is
which prime factors can cancel. The exact identity in
Lemma~\ref{lem:denominator-amplification} controls that cancellation;
Theorem~\ref{thm:denominator-period-bound} then propagates the resulting
constraints through strictly decreasing primes. This gives an explicit arithmetic obstruction and applies it to the
historical question.
The tools---geometric series, prime valuations, and multiplicative
orders---are classical; no new foundational principle is claimed.

The research began with the weighted-tail viewpoint in manuscript~[090]
of the supplied corpus~\cite{Corpus090}, inspected at the fixed commit
recorded in the bibliography. That corpus is used as a source of research
directions, not as an assumption of correctness. Every assertion needed
for the present result is proved below.


\section{Denominators and eventual periods in weighted binary expansions}
\label{sec:weighted-periods}

For a binary sequence $\varepsilon=(\varepsilon_n)_{n\ge1}$, write
\[
 X(\varepsilon)=\sum_{n\ge1}\frac{n\varepsilon_n}{2^n}.
\]
Every $x\in[0,2]$ admits such a representation; the elementary greedy construction is recalled below. We determine an explicit restriction on every eventual period of a representation of a fixed rational number. In particular, $1/3$ has no eventually periodic representation. This answers the specific periodicity possibility discussed by Borwein and Loring on printed page~391 of~\cite{BL1990}. It does not settle either of the remaining questions in Erd\H{o}s Problem~\#261: universal finite representability of $n2^{-n}$ or continuum multiplicity for some rational value.

The periodic-tail formula used below is classical generating-function algebra. In the closely related recent literature, Wang and Grau Ribas~\cite[Lemma~A.12 and Propositions~7.2/A.15]{Wang2026v1} use this formula and the identification of the primitive binary period with a multiplicative order to obtain a density estimate. The argument here extracts prime valuations from the same formula and propagates them through the prime divisors of multiplicative orders. No claim that the underlying generating-function identity is new is intended.

For an odd prime $p$, define
\[
 t_p=\operatorname{ord}_p(2),\qquad
 w_p=v_p(2^{t_p}-1).
\]
Here $v_p$ is the usual additive valuation, with $v_p(0)=+\infty$. In particular $w_p\ge1$ and $t_p\mid p-1$.

\begin{lemma}[Periodic tail formula]\label{lem:periodic-tail-formula}
Suppose that $\varepsilon_{N+i}=e_i$ for $i\ge1$, where $(e_i)$ is periodic of minimal period $r$, with digits in $\{0,1\}$. Put
\[
 M=2^r-1,\qquad
 S=\sum_{i=1}^r e_i2^{r-i},\qquad
 T=\sum_{i=1}^r i e_i2^{r-i},
\]
\[
 y=\sum_{i\ge1}\frac{e_i}{2^i}=\frac{S}{M},\qquad
 z=\sum_{i\ge1}\frac{i e_i}{2^i}
   =\frac{MT+rS}{M^2}.
\]
Then
\begin{equation}\label{eq:periodic-tail-shift}
 2^NX(\varepsilon)=I+Ny+z,
 \qquad I=\sum_{n=1}^{N}n\varepsilon_n2^{N-n}\in\mathbb Z.
\end{equation}
If the periodic word is nonconstant and $d$ is the reduced denominator of $y$, then $d>1$ is odd and
\[
 r=\operatorname{ord}_d(2)<d.
\]
\end{lemma}

\begin{proof}
Summing over positions $i+kr$ gives
\[
 y=\sum_{i=1}^r\frac{e_i2^{-i}}{1-2^{-r}},\qquad
 z=\sum_{i=1}^r e_i2^{-i}
 \left(\frac{i}{1-2^{-r}}+
       \frac{r2^{-r}}{(1-2^{-r})^2}\right),
\]
which yields the displayed formulas. Separating the finite prefix proves~\eqref{eq:periodic-tail-shift}. A nonconstant periodic binary word represents a rational $y$ strictly between zero and one. Multiplication by $2^s$ shifts the ordinary binary digits by $s$ places; its fractional part returns to $y$ precisely when $d\mid 2^s-1$. Uniqueness of the nonterminating binary expansion therefore identifies its minimal period with the multiplicative order of $2$ modulo $d$. Euler's theorem gives $r\le\varphi(d)<d$.
\end{proof}


\begin{lemma}[Order lifting]\label{lem:order-lifting}
For an odd prime $p$ and every $a\ge1$,
\[
 \ord_{p^a}(2)=t_p p^{\max(0,a-w_p)}.
\]
Moreover, if $t_p\mid r$, then
$v_p(2^r-1)=w_p+v_p(r)$.
\end{lemma}
\begin{proof}
For $u\equiv1\pmod p$ and odd $p$, the binomial theorem gives
$v_p(u^m-1)=v_p(u-1)+v_p(m)$. To see this, an exponent coprime to
$p$ leaves the valuation unchanged, while raising to the power $p$
increases it by exactly one; iteration proves the identity.
Apply it to $u=2^{t_p}$. Any exponent giving $1$ modulo $p^a$ is a
multiple $t_pm$, and it works exactly when $w_p+v_p(m)\ge a$.
The least such multiple proves the order formula. The final assertion
also uses $p\nmid t_p$.
\end{proof}

\begin{lemma}[Exact denominator amplification]\label{lem:denominator-amplification}
Under the nonconstant-tail hypotheses of Lemma~\ref{lem:periodic-tail-formula}, let $Q$ be the odd part of the reduced denominator of $X(\varepsilon)$. For an odd prime $p\mid d$, put
\[
 a=v_p(d),\qquad b=v_p(r).
\]
If $a>b$, then
\begin{equation}\label{eq:exact-amplification}
 v_p(Q)=a+w_p.
\end{equation}
In particular, the odd denominator of $X(\varepsilon)$ is divisible by the square of some prime.
\end{lemma}

\begin{proof}
Write $k=v_p(M)$. Reduction of $S/M$ gives $v_p(S)=k-a$. Since $a>b$,
\[
 v_p(rS)=k-a+b<k\le v_p(MT).
\]
Thus the two numerator summands in $MT+rS$ cannot cancel at their least valuation, and
\[
 v_p(z)=k-a+b-2k=-k-a+b.
\]
As $p\mid d$, the number $t_p$ divides $r$. Since $p\nmid t_p$, the lifting-the-exponent identity gives
\[
 k=v_p(2^r-1)=w_p+v_p(r/t_p)=w_p+b.
\]
Consequently $v_p(z)=-a-w_p$. On the other hand $v_p(Ny)\ge-a$ and $v_p(I)\ge0$. These terms cannot cancel the strictly smaller valuation of $z$ in~\eqref{eq:periodic-tail-shift}. Multiplication by $2^N$ leaves every odd-prime valuation unchanged, proving~\eqref{eq:exact-amplification}.

Finally, $r<d$ implies $d\nmid r$. Hence there is a prime $p\mid d$ with $v_p(d)>v_p(r)$. For this prime, $v_p(Q)=a+w_p\ge2$.
\end{proof}

\begin{corollary}[Squarefree-denominator obstruction]\label{cor:squarefree-periodicity}
Let $x\in\mathbb Q\cap[0,2]$, and suppose that the odd part $Q>1$ of its reduced denominator is squarefree. Then no binary representation
\[
 x=\sum_{n\ge1}n\varepsilon_n2^{-n}
\]
is eventually periodic. In particular, every such representation of $1/3$ is aperiodic after every finite prefix.
\end{corollary}

\begin{proof}
A nonconstant periodic tail contradicts Lemma~\ref{lem:denominator-amplification}. An eventually constant binary sequence gives a dyadic value: a zero tail is finite, while a one tail contributes
\[
 \sum_{n>N}\frac n{2^n}=\frac{N+2}{2^N}.
\]
Either is incompatible with $Q>1$.
\end{proof}

For completeness, existence of the representations appearing in the corollary follows directly from a tail interval. Start with $\rho_0=x$ and at step $n\ge1$ choose a digit $\varepsilon_n$ for which
\[
 \rho_n=2\rho_{n-1}-n\varepsilon_n\in[0,n+2].
\]
If $\rho_{n-1}\in[0,n+1]$, exclusion is possible for $2\rho_{n-1}\le n+2$ and inclusion is possible for $2\rho_{n-1}\ge n$; these conditions cover the interval. Then
\[
 x-\sum_{j=1}^n\frac{j\varepsilon_j}{2^j}
 =2^{-n}\rho_n\longrightarrow0.
\]
Thus the conclusion for $1/3$ is not vacuous.

\begin{definition}[Prime set and period bound attached to a denominator]
\label{def:denominator-period-bound}
For an odd positive integer $Q$, let $\mathcal P_Q$ be the smallest set of odd primes with the following properties:
\begin{enumerate}
\item it contains every prime $p$ with $v_p(Q)>w_p$;
\item if $q\in\mathcal P_Q$, it contains every odd prime divisor of $t_q$.
\end{enumerate}
This set is finite, since $t_q\mid q-1$ and hence every newly added prime is smaller than its parent. For $p\in\mathcal P_Q$ set
\[
 e_p=\max\left\{0,\ v_p(Q)-w_p,\
                 \max_{\substack{q\in\mathcal P_Q\\q>p}}v_p(t_q)\right\},
 \qquad D_Q=\prod_{p\in\mathcal P_Q}p^{e_p}.
\]
An empty inner maximum is zero and an empty product is one. Every prime in $\mathcal P_Q$ has $e_p\ge1$. If $D_Q>1$, put $R_Q=\operatorname{ord}_{D_Q}(2)$.
\end{definition}

\begin{theorem}[A bound for every eventual period]\label{thm:denominator-period-bound}
Let $x$ be rational with reduced denominator $2^sQ$, where $Q$ is odd. Suppose a binary representation of $x$ has a nonconstant eventual period. Let $d$ be the reduced denominator of the ordinary binary value of a periodic tail, and let $r$ be its minimal period. Then
\[
 d\mid D_Q.
\]
Consequently $D_Q>1$ and
\[
 r\mid R_Q.
\]
The bound depends only on the odd denominator of $x$, independently of the numerator, the dyadic denominator, the preperiod, and the chosen representation.
\end{theorem}

\begin{proof}
For each prime $p\mid d$, put $a_p=v_p(d)$ and $b_p=v_p(r)$. We first prove the following alternative:
\begin{equation}\label{eq:prime-alternative}
 a_p\le v_p(Q)-w_p
 \quad\hbox{or}\quad
 a_p\le v_p(t_q)\text{ for some prime }q\mid d\text{ with }q>p.
\end{equation}
If $a_p>b_p$, Lemma~\ref{lem:denominator-amplification} gives equality in the first alternative. Otherwise $a_p\le b_p$. The standard order formula at odd prime powers is
\[
 \operatorname{ord}_{q^a}(2)
   =t_q q^{\max\{0,a-w_q\}}.
\]
Since $r=\operatorname{lcm}_{q\mid d}\operatorname{ord}_{q^{a_q}}(2)$, taking its $p$-valuation yields
\[
 b_p=\max\left\{\max\{0,a_p-w_p\},\
                \max_{\substack{q\mid d\\q>p}}v_p(t_q)\right\}.
\]
Here $q<p$ contributes zero because $t_q\mid q-1<p$, and $t_p$ is coprime to $p$. The first entry is strictly smaller than $a_p$. Thus $b_p\ge a_p$ forces the second alternative in~\eqref{eq:prime-alternative}.

Now process the prime divisors of $d$ in decreasing order. The largest one must satisfy the first alternative, so it is a seed of $\mathcal P_Q$ and its exponent is at most $e_p$. For any smaller prime $p$, the first alternative again makes it a seed with the required exponent bound. In the second alternative, the larger prime $q$ has already been placed in $\mathcal P_Q$; since $p\mid t_q$, the closure rule places $p$ there too, and the definition of $e_p$ bounds $a_p$. This proves $d\mid D_Q$. Multiplicative orders under divisibility then give
\[
 r=\operatorname{ord}_d(2)\mid\operatorname{ord}_{D_Q}(2)=R_Q.
\]
\end{proof}

\begin{corollary}[A finite list of possible tail words]\label{cor:finite-tail-alphabet}
For a fixed odd denominator $Q$, every ordinary binary value of a nonconstant periodic tail belongs to the finite set
\[
 \left\{\frac{j}{D_Q}:1\le j<D_Q\right\}.
\]
Thus finitely many periodic binary tail sequences can occur, including their phases. This does not bound the point at which a tail begins.
\end{corollary}

\begin{proof}
The reduced denominator $d$ of the tail value divides $D_Q$, and the value lies strictly between zero and one. Each such rational has a unique nonterminating, purely periodic binary expansion.
\end{proof}

\begin{example}[The squarefree obstruction has a real boundary]
Although $1/3$ has no eventually periodic representation, $1/9$ does:
\[
 \frac19=\frac7{2^7}+\frac8{2^8}
          +\sum_{k\ge0}\frac{9+2k}{2^{9+2k}}.
\]
Indeed the finite prefix contributes $11/128$. The tail begins after $N=8$ and has periodic word $10$, for which $y=2/3$ and $z=10/9$. Its weighted contribution is therefore
\[
 2^{-8}(8y+z)=\frac{29}{1152},\qquad
 \frac{11}{128}+\frac{29}{1152}=\frac19.
\]
\end{example}

\begin{remark}[Scope and examples]
Theorem~\ref{thm:denominator-period-bound} is a necessary condition for periodic representability, not a criterion asserting its existence for each rational numerator. It bounds the period and not the preperiod. For example,
\[
\begin{array}{c|ccccc}
 Q&9&25&49&81&121\\\hline
 D_Q&3&5&21&27&55\\
 R_Q&2&4&6&18&20
\end{array}
\]
where the primes occurring here have $w_p=1$. Thus every nonconstant eventual period at odd denominator $49$ divides $6$, and at odd denominator $121$ it divides $20$. Squarefree $Q$ has no seed, giving $D_Q=1$ and recovering the obstruction above. These conclusions concern a classical periodicity question adjacent to Erd\H{o}s~\#261; they provide no proof of infinite branching or continuum multiplicity.
\end{remark}

\section{Examples and the role of higher prime valuations}
The period bound is attained for some denominators. The purely periodic
word $0011$ has $r=4$, $M=15$, $S=3$, and $T=10$, so its weighted
value is $18/25$. Its minimal period is $4=R_{25}$. For $Q=49$, take
a zero prefix of length one followed by the repeated word $000011$.
Here $y=1/21$ and $z=38/147$, and hence
\[
 x=\frac12\left(\frac1{21}+\frac{38}{147}\right)=\frac{15}{98}.
\]
Its minimal period is $6=R_{49}$. The prime $3$ in $D_{49}=21$ comes
from $\ord_7(2)=3$; it need not divide $Q$ itself. Thus simply retaining
the primes dividing $Q$ would lose valid tails.

The use of $w_p$, rather than the weaker uniform lower bound $w_p\ge1$,
also matters. For the prime $p=1093$, exact modular arithmetic gives
\[
\begin{array}{c|rrrr}
 m&364&182&52&28\\\hline
 2^m\bmod1093&1&1092&27&121
\end{array}
\]
and
\[
 2^{364}\equiv1\pmod{1093^2},\qquad
 2^{364}\equiv1064432260\not\equiv1\pmod{1093^3}.
\]
Trial division proves that $1093$ is prime. Since the prime divisors of
$364$ are $2,7,13$, the table proves $t_{1093}=364$, and the second
line proves $w_{1093}=2$. For $Q=1093^2$ there is consequently no seed
in Definition~\ref{def:denominator-period-bound}; thus $D_Q=1$.

\begin{corollary}
No weighted binary representation of $1/1093^2$ is eventually periodic.
\end{corollary}
\begin{proof}
Theorem~\ref{thm:denominator-period-bound} excludes a nonconstant
periodic tail, and a constant tail would give a dyadic rational.
\end{proof}
This conclusion is not covered by the squarefree-denominator corollary.
The number-theoretic property of $1093$ is used only as an explicitly
checked example, not as an unproved input.


\section{The denominator obstruction in every integer base}
\label{sec:general-base}

Fix an integer $b\geq 2$ and use the full digit alphabet
$\{0,1,\ldots,b-1\}$.  A weighted base-$b$ representation is an expansion
\[
 x=\sum_{n=1}^{\infty}\frac{n\varepsilon_n}{b^n},
 \qquad \varepsilon_n\in\{0,1,\ldots,b-1\}.
\]
If $x$ is rational with reduced positive denominator $q$, write
\[
 Q=\prod_{\substack{p\mid q\\p\nmid b}}p^{v_p(q)},
\]
where $v_p$ denotes the usual prime-adic valuation.  Thus $Q$ is the part
of the denominator coprime to $b$.

\begin{theorem}\label{thm:general-base-obstruction}
Let $x$ have an eventually periodic weighted base-$b$ representation.
If its eventual tail is neither constantly $0$ nor constantly $b-1$,
then the reduced denominator of $x$ is divisible by $p^2$ for some prime
$p\nmid b$.  If the eventual tail is constantly $0$ or constantly $b-1$,
then $Q\mid b-1$.

Consequently, if a rational number $x$ has squarefree $Q$ and
$Q\nmid b-1$, then none of its weighted base-$b$ representations is
eventually periodic.
\end{theorem}

\begin{proof}
Choose $N\geq0$ such that $\delta_i=\varepsilon_{N+i}$ is purely
periodic, and let $r$ be its least positive period.  Put
\[
 M=b^r-1,\qquad
 S=\sum_{i=1}^{r}\delta_i b^{r-i},\qquad
 T=\sum_{i=1}^{r}i\delta_i b^{r-i}.
\]
Summing the geometric series and its derivative gives
\begin{equation}\label{eq:general-base-periodic-tail}
 y:=\sum_{i=1}^{\infty}\frac{\delta_i}{b^i}=\frac{S}{M},
 \qquad
 z:=\sum_{i=1}^{\infty}\frac{i\delta_i}{b^i}
     =\frac{MT+rS}{M^2}.
\end{equation}
Writing the finite prefix as
$F=\sum_{n=1}^{N}n\varepsilon_n b^{-n}$, we therefore have
\begin{equation}\label{eq:general-base-shift}
 x=F+b^{-N}(z+Ny).
\end{equation}

First suppose that the periodic tail is neither all $0$ nor all $b-1$.
Then $0<y<1$.  Its reduced denominator $d$ divides $M$ and satisfies
$d>1$, so $\gcd(d,b)=1$.  The least period of the ordinary base-$b$
expansion of $y$ is
\[
 r=\operatorname{ord}_d(b)\leq\varphi(d)<d.
\]
Indeed, $b^s y-y$ is an integer precisely when $d\mid b^s-1$;
uniqueness of the ordinary base-$b$ expansion of this nonterminating
rational identifies these integers $s$ with its periods.  In particular,
$d\nmid r$.  We may therefore choose a prime $p\mid d$ such that
\[
 a:=v_p(d)>h:=v_p(r).
\]
Let $k=v_p(M)$.  Reduction of $S/M$ gives $v_p(S)=k-a$, and hence
\[
 v_p(rS)=k-a+h<k\leq v_p(MT).
\]
The two summands in the numerator of $z$ have distinct $p$-adic
valuations.  Thus
\[
 v_p(MT+rS)=k-a+h,
\]
and the exponent of $p$ in the reduced denominator of $z$ is exactly
\[
 E=2k-(k-a+h)=k+a-h>a\geq1.
\]
The denominator of $Ny$ has $p$-exponent at most $a$, so addition of
$Ny$ cannot cancel the larger exponent $E$ in the denominator of $z$.
Moreover, $p\nmid b$, since $p\mid M$.  Multiplication by $b^{-N}$ and
addition of the finite base-$b$ rational $F$ likewise leave this exponent
unchanged.  Equation~\eqref{eq:general-base-shift} therefore proves that
$p^2$ divides the reduced denominator of $x$.

If the tail is all $0$, then $x=F$ and $Q=1$.  If the tail is all
$b-1$, direct summation gives
\[
 y=1,\qquad z=\frac{b}{b-1},\qquad
 x=F+b^{-N}\left(N+\frac{b}{b-1}\right).
\]
All denominator factors outside the primes dividing $b$ consequently
divide $b-1$.  This proves the two exceptional cases and the stated
consequence.
\end{proof}

The extremal-tail exception is necessary.  For example, in base $3$
the constant digit sequence $\varepsilon_n=2$ represents
$\sum_{n\geq1}2n3^{-n}=3/2$, whose reduced denominator is squarefree.
Constant interior digits $1,\ldots,b-2$, when they exist, are covered
by the nonextremal case of Theorem~\ref{thm:general-base-obstruction}.

\begin{lemma}\label{lem:general-base-existence}
Every real number $x\in[0,b/(b-1)]$ has a weighted base-$b$
representation with digits in $\{0,1,\ldots,b-1\}$.
\end{lemma}

\begin{proof}
Set $w_n=n/b^n$ and let
\[
 H_n=(b-1)\sum_{j>n}\frac{j}{b^j}
     =b^{-n}\left(n+\frac{b}{b-1}\right)
     \qquad(n\geq0).
\]
Thus $H_0=b/(b-1)$, $H_n\to0$, and
$H_n>w_n$ for $n\geq1$.  The intervals
\[
 [d w_n,d w_n+H_n],\qquad d=0,1,\ldots,b-1,
\]
overlap and cover
$[0,(b-1)w_n+H_n]=[0,H_{n-1}]$.
Starting with $R_0=x$, choose at step $n$ a digit $\varepsilon_n$
whose interval contains $R_{n-1}$, and put
$R_n=R_{n-1}-\varepsilon_n w_n$.  Then $0\leq R_n\leq H_n$.
For example, the explicit greedy choice
\[
 \varepsilon_n=
 \min\left\{b-1,\left\lfloor\frac{R_{n-1}}{w_n}\right\rfloor\right\}
\]
has this property.  Since
$x-\sum_{j=1}^{n}\varepsilon_j w_j=R_n\to0$, the desired expansion
follows.
\end{proof}

In particular, for every prime $p\nmid b(b-1)$ the number $1/p$ lies
in the interval of Lemma~\ref{lem:general-base-existence} and has at
least one weighted base-$b$ representation, but none is eventually
periodic.  In base $2$ one may take $p=3$.

\section{Periodic representations are isolated}
\label{sec:periodic-isolation}

For $x\in[0,2]$, equip the representation fiber
\[
 \mathcal F_x=
 \left\{(\varepsilon_n)\in\{0,1\}^{\mathbb N}:
       x=\sum_{n=1}^{\infty}\frac{n\varepsilon_n}{2^n}\right\}
\]
with the subspace topology inherited from the product topology on
$\{0,1\}^{\mathbb N}$.  A representation is isolated if some finite
initial digit block distinguishes it from every other member of
$\mathcal F_x$.

\begin{theorem}\label{thm:periodic-isolation}
Every eventually periodic binary weighted representation is an
isolated point of its full representation fiber.
\end{theorem}

\begin{proof}
Let $\varepsilon\in\mathcal F_x$ be eventually periodic.  If its tail
is constantly $0$, then fixing the digits through the start of that
tail leaves the smallest possible remaining sum, namely zero.
Every subsequent digit is therefore forced.  The same argument with
the largest possible remaining sum applies to a tail constantly $1$.
Both cases give isolation.

Otherwise choose $N\geq0$ and $r\geq1$ such that the tail after $N$
is periodic with period $r$ and contains both digits.  Choose an integer
$L\geq\max\{N,2^{r+1}\}$.  Suppose that another representation
$\eta\in\mathcal F_x$ agrees with $\varepsilon$ through index $L$,
and let $n>L$ be their first differing index.

If $\varepsilon_n=1$ and $\eta_n=0$, equality of the represented
numbers would require the loss $n2^{-n}$ to be compensated by later
changes.  The largest available compensation is obtained by turning
every later zero of $\varepsilon$ into a one.  Every block of $r$
consecutive digits in this tail contains a one, so in particular
\[
 \sum_{t\geq1}\frac{1-\varepsilon_{n+t}}{2^t}
 \leq 1-2^{-r}.
\]
Using also $\sum_{t\geq1}t2^{-t}=2$, we obtain
\[
 \begin{split}
 2^n\sum_{m>n}\frac{m(1-\varepsilon_m)}{2^m}
 &=\sum_{t\geq1}\frac{(n+t)(1-\varepsilon_{n+t})}{2^t}\\
 &\leq n(1-2^{-r})+2<n,
 \end{split}
\]
where the strict inequality follows from $n>2^{r+1}$.  Compensation
is impossible.

If $\varepsilon_n=0$ and $\eta_n=1$, the required compensation would
instead come from deleting later ones.  Since every block of $r$
tail digits contains a zero, the identical estimate with
$\varepsilon_{n+t}$ in place of $1-\varepsilon_{n+t}$ again bounds
the available compensation strictly below $n2^{-n}$.
Both cases are impossible.  Thus the cylinder fixing the first $L$
digits meets $\mathcal F_x$ only at $\varepsilon$, proving isolation.
\end{proof}

This theorem gives only one direction: eventual periodicity implies
isolation.  It does not prove that an isolated representation of a
rational number is eventually periodic.  In particular, excluding
eventually periodic representations of $1/3$ does not establish that
its representation fiber is perfect or uncountable.

\section{Literature, scope, and verification}
\subsection{What has and has not been established}
The completed mathematical result is an obstruction to eventual
periodicity for all representations of the stated rational numbers,
together with a denominator-dependent bound on every possible
nonconstant period. The original universal periodicity possibility
is false. The statements are unconditional and their proofs do not
depend on the finite calculations.

Neither remaining assertion of Erd\H{o}s~\#261 is proved here. In the
formal-conjectures snapshot~\cite{Formal261}, these ask whether every
positive integer $n$ has a nontrivial finite representation of $n2^{-n}$,
and whether some rational number has continuum many representations
by increasing positive integers. Aperiodicity is a property of one
sequence; it does not imply that a value has many representations.
The missing converse discussed in Section~\ref{sec:periodic-isolation}
must not be silently assumed. We therefore do not present this work
as fulfilling the requested conjunction of a still-open named
Erd\H{o}s problem, a complete solution, and certified historical priority.

\subsection{Relation to previous methods}
The closest identified antecedent is Wang and Grau Ribas~\cite{Wang2026v1},
specifically Lemma~A.12 and Proposition~A.15 in version~1. It contains
the periodic-tail identity and multiplicative-order reduction used in
Section~\ref{sec:weighted-periods}. Those ingredients are credited as
prior work. The additional conclusions developed here are the exact
valuation obstruction and its finite propagation through prime
divisors of orders. We use no density theorem from that preprint.

The literature check examined Borwein--Loring, the 2020 paper of
Tengely--Ulas--Zygad\l{}o~\cite{TUZ2020}, and all four available versions
of arXiv:2606.24972, together with title, citation, and phrase searches
for weighted or $*$-binary expansions, eventual periodicity, $1/3$,
and squarefree denominators. No earlier identical resolution or
period-divisor theorem was found in those sources. This is a bounded
search result, not a proof that no prior solution exists. The original
source documents that the periodicity possibility was unresolved there
in 1990; this manuscript does not certify that it remained unresolved
through every intervening publication. In particular, the argument
does not establish a historically new foundational principle.

\subsection{Exact checks and reproducibility}
Appendix~\ref{app:code} contains the full standard-library Python
verifier. It enumerates every nonconstant binary word of lengths
$2$ through $17$, including words with a smaller minimal period.
For lengths at most $8$, every shift $N=0,\ldots,2^r-2$ is tested;
for larger lengths it uses
$N\in\{0,1,2,r,r+1,2^r-2,1000007\}$.
It computes the reduced odd denominator exactly, checks $d\mid D_Q$
and $r_{\min}\mid R_Q$, and checks the exact valuation identity
whenever its hypothesis holds. The completed run produced
\[
\begin{array}{lr}
\toprule
\text{Nonconstant words}&262108\\
\text{Word--shift cases}&1917164\\
\text{Exact valuation checks}&3810912\\
\text{Failed assertions}&0\\
\bottomrule
\end{array}
\]
Adding a finite prefix contributes an integer to the scaled expression
and cannot change its reduced odd denominator; the dyadic scaling
likewise leaves that denominator unchanged. The checker therefore
tests the relevant arithmetic directly. The finite counts are
consistency checks only. All word lengths, shifts, and prefixes are
covered by the proofs, not by an extrapolation of these counts.

\appendix
\section{Self-contained exact verifier}\label{app:code}
Save the following listing as \texttt{verify\_period\_bounds.py} and run with Python~3.9 or later:
\begin{center}\texttt{python verify\_period\_bounds.py 17}\end{center}
It uses only integer arithmetic. The default argument is $17$; a
smaller argument shortens the diagnostic run. The verifier is embedded
in this single \LaTeX{} source and needs no external data.
% VERIFIER-BEGIN
\begin{lstlisting}

"""Independent exact checks for the proposed weighted-period bound.

No numerical fitting or floating point is used. This tests finite words; the
mathematical proof is the separate descending-prime/valuation argument.
"""
from functools import lru_cache
from math import gcd, lcm
import json
import sys


@lru_cache(None)
def factor(n):
    out = []
    p = 2
    while p * p <= n:
        a = 0
        while n % p == 0:
            n //= p
            a += 1
        if a:
            out.append((p, a))
        p = 3 if p == 2 else p + 2
    if n > 1:
        out.append((n, 1))
    return tuple(out)


def vp(n, p):
    assert n > 0
    a = 0
    while n % p == 0:
        n //= p
        a += 1
    return a


@lru_cache(None)
def prime_data(p):
    assert p > 2
    t = p - 1
    for q, _ in factor(p - 1):
        while t % q == 0 and pow(2, t // q, p) == 1:
            t //= q
    w = 1
    while pow(2, t, p ** (w + 1)) == 1:
        w += 1
    return t, w


@lru_cache(None)
def bound(Q):
    Qf = dict(factor(Q))
    seeds = {p for p, a in Qf.items() if a > prime_data(p)[1]}
    closure = set(seeds)
    pending = list(seeds)
    while pending:
        q = pending.pop()
        tq, _ = prime_data(q)
        for p, _ in factor(tq):
            if p != 2 and p not in closure:
                closure.add(p)
                pending.append(p)
    D = 1
    period = 1
    exponents = {}
    for p in closure:
        tp, wp = prime_data(p)
        incoming = [vp(prime_data(q)[0], p) for q in closure if q > p]
        ep = max([Qf.get(p, 0) - wp, 0] + incoming)
        assert ep > 0
        exponents[p] = ep
        D *= p ** ep
        period = lcm(period, tp * p ** max(0, ep - wp))
    return D, period, exponents


def run(max_period=17):
    cases = 0
    words = 0
    exact_valuation_checks = 0
    for r in range(2, max_period + 1):
        M = 2 ** r - 1
        for S in range(1, M):
            words += 1
            d = M // gcd(S, M)
            df = dict(factor(d))
            minimal_r = 1
            for p, a in df.items():
                tp, wp = prime_data(p)
                minimal_r = lcm(minimal_r, tp * p ** max(0, a - wp))
            T = sum(i * 2 ** (r - i) for i in range(1, r + 1)
                    if (S >> (r - i)) & 1)
            # Exhaust every shift residue for r<=8. For larger periods use
            # several shifts, including a much larger, irregular one.
            shifts = range(M) if r <= 8 else (0, 1, 2, r, r + 1, M - 1, 10**6 + 7)
            for N in shifts:
                num = M * (T + N * S) + r * S
                Q = M * M // gcd(M * M, num)
                D, period, _ = bound(Q)
                assert D % d == 0, (r, S, N, d, Q, D)
                assert period % minimal_r == 0, (r, S, N, minimal_r, Q, period)
                for p, a in df.items():
                    b = vp(minimal_r, p)
                    if a > b:
                        _, w = prime_data(p)
                        assert vp(Q, p) == a + w, (r, S, N, p, Q, a, b, w)
                        exact_valuation_checks += 1
                cases += 1
        print(f"r={r}: cumulative cases={cases}", flush=True)
    # An important boundary example: seed 7 must bring in prime 3.
    assert bound(49)[:2] == (21, 6)
    # Explicit Wieferich-prime cases exercise w_p>1 in the construction.
    assert prime_data(1093)[1] >= 2
    assert bound(1093**2)[0] == 1
    assert bound(1093**3)[0] % 1093 == 0
    report = {
        "max_word_length": max_period,
        "nonconstant_words": words,
        "word_shift_cases": cases,
        "exact_valuation_checks": exact_valuation_checks,
        "failures": 0,
        "example_Q49": {"D": bound(49)[0], "period_bound": bound(49)[1]},
        "wieferich_1093": prime_data(1093),
    }
    print(json.dumps(report, indent=2))


if __name__ == '__main__':
    run(int(sys.argv[1]) if len(sys.argv) > 1 else 17)

\end{lstlisting}
% VERIFIER-END

\begin{thebibliography}{9}
\raggedright
\bibitem{BL1990}
P.~B. Borwein and T.~A. Loring,
\emph{Some questions of Erd\H{o}s and Graham on numbers of the form
$\sum g_n/2^{g_n}$},
Mathematics of Computation \textbf{54} (189) (1990), 377--394.
\href{https://doi.org/10.1090/S0025-5718-1990-0990598-9}{doi:10.1090/S0025-5718-1990-0990598-9}.
\href{https://www.cecm.sfu.ca/~pborwein/PAPERS/P46.pdf}{Author-hosted PDF}.

\bibitem{Wang2026v1}
H. Wang and J.~M. Grau Ribas,
\emph{Positive dyadic density for rational weighted binary expansions},
arXiv:2606.24972v1 [math.NT], 23 June 2026;
especially Lemma~A.12 (p.~21) and Proposition~A.15 (p.~22).
\url{https://arxiv.org/pdf/2606.24972v1}.
Version~1 is specified because later versions change the title,
authorship, and organization.

\bibitem{TUZ2020}
S. Tengely, M. Ulas, and J. Zygad\l{}o,
\emph{On a Diophantine equation of Erd\H{o}s and Graham},
Journal of Number Theory \textbf{217} (2020), 445--459.
\href{https://doi.org/10.1016/j.jnt.2020.05.006}{doi:10.1016/j.jnt.2020.05.006}.
\url{https://arxiv.org/abs/2008.01501}.

\bibitem{Corpus090}
Keio, \emph{Tent rigidity for the Erd\H{o}s--Graham equation
$n/2^n=\sum_i a_i/2^{a_i}$ and the Tengely--Ulas--Zygad\l{}o conjecture},
29 October 2026, manuscript~[090] in
\emph{Verifiable Mathematical Reasoning Corpus}.
Repository snapshot inspected 2 November 2026:
\href{https://github.com/QEEDDD/Verifiable-Mathematical-Reasoning-Corpus/tree/ecd87de3348dcb706a76bb12fde722d7a348a012}{commit \texttt{ecd87de3348dcb706a76bb12fde722d7a348a012}}.
This is a research-direction source, not a theorem dependency.

\bibitem{Formal261}
Google DeepMind and contributors, \emph{Formal Conjectures},
\nolinkurl{FormalConjectures/ErdosProblems/261.lean},
snapshot inspected 2 November 2026,
\href{https://github.com/google-deepmind/formal-conjectures/blob/d995ef5d306637c0a2e489c7e579647168d7df2f/FormalConjectures/ErdosProblems/261.lean}{commit \texttt{d995ef5d306637c0a2e489c7e579647168d7df2f}}.
Used for the precise separation of the numbered problem's parts;
not a proof source for the present theorems.
\end{thebibliography}
\end{document}
